In what follows, let $(\mx,\mathcal{A})$ be a measurable space a consider a family $\big(P_\theta\big)_{\theta\in\Theta}$ of probability measures on the latter measurable space. Let us have an observation $X\in \mx$ generated by $P_\theta$ for some unknown parameter $\theta \in \Theta$. 
\section{Gaussian sequence model} 
\begin{definition}
    We observe a sequence $X = (X_i)_{i\in I} \in \ell^2$ which is of the following form : 
    \begin{equation}\label{gaussmodel}
        X = v + \eta
    \end{equation}
    Here, $v\in \ell^2$ is the (unknown) mean vector, $\eta_i \stackrel{i.i.d}{\sim} \mathcal{N}(0,1)$ 
\end{definition}
\noindent
We denote by $P_v$ the distribution of $X$ when the signal is $v$, and by $P_0$ the distribution of $X$ when $v = 0$. Naturally, one can define the following sets : 
\begin{equation*}
    \mpg_n \coloneqq \bigl\{P_v, v \in \R^n\bigr\} \text{ and } \mpg_\infty \coloneqq \bigl\{P_v, v\in \ell^2\bigr\}
\end{equation*}

\begin{remark}
    Note that in the case where $I$ is countably infinite, we often assume that $I = \N$, and $\ell^2$ is the Hilbert space of square integrable sequences i.e. : 
\begin{equation*}
    \ell^2 \coloneqq \Bigl\{ (u_i)_{i\in I} : \sum_{i \in I} u_i^2 < \infty \Bigr\}
\end{equation*}
\end{remark}
\subsection{The Sobolev ellipsoid and the Gaussian sequence model}

\begin{definition}[Gaussian sequence model on the Sobolev ellipsoid]
For $s, C > 0$, the \emph{Sobolev ellipsoid} is $\mathcal{E}(s,C) := \bigl\{\theta \in
\mathbb{R}^{\mathbb{N}} : \sum_{j=1}^\infty j^{2s}\theta_j^2 \leq C\bigr\}$. Setting
$\mu_\theta = \bigotimes_{j=1}^\infty \mathcal{N}(\theta_j, 1)$, the model class is:
\[
    \mathcal{P}_G(s, C_G) := \{\mu_\theta \;;\; \theta \in \mathcal{E}(s, C_G)\}.
\]
\end{definition}

\noindent\textbf{Intuition.} The weight $j^{2s}$ in the Sobolev constraint penalizes
high-frequency components: for the sum to stay bounded, the coefficients must satisfy
$\theta_j = O(j^{-s})$. A large $s$ forces fast decay (smooth signal, easy to detect),
while a small $s$ allows more energy at high frequencies (rough signal, harder to detect).
This is directly reflected in Theorem~\ref{thm:main}: the exponent $\frac{4s+1}{2s}$ is
decreasing in $s$.

\medskip\noindent\textbf{Connection to the white noise model.}
This model arises naturally from the Gaussian white noise model $dY_t = f(t)\,dt + dW_t$.
Projecting onto an orthonormal basis $\{\varphi_j\}$ of $L^2[0,1]$ gives
$y_j = \theta_j + \varepsilon_j$ with $\varepsilon_j \overset{\text{i.i.d.}}{\sim}
\mathcal{N}(0,1)$ and $\theta_j = \langle f, \varphi_j\rangle$. The Sobolev constraint on
$\theta$ then corresponds exactly to a Sobolev regularity condition on $f$. Testing $H_0 : P_X = P_0$ is thus equivalent to testing whether an unknown signal $f$
equals a given reference $f_0$.

\medskip\noindent\textbf{Truncation.} Since $\theta \in \mathcal{E}(s,C)$ forces
$\theta_j \to 0$, only the first $J \asymp \varepsilon^{-1/s}$ coordinates carry a detectable
signal; beyond index $J$ the components are swamped by noise. This truncation is the key
structural insight behind the separating set constructed in what follows.

\section{Classical theory of minimax distinguishability}
\subsection{Asymptotic Framework}
In this chapter, instead of considering only one hypothesis testing problem, we will follow an asymptotic approach in which we will be considering a family of Hypothesis testing problems indexed by a parameter $\epsilon \rightarrow \epsilon_0$. 
Depending on the setting, this parameter may be the dimensions $n \rightarrow \infty$, the sample size $N \rightarrow \infty$ or the white noise level $\epsilon \rightarrow 0$. More generally, this corresponds to increasing the amount of available information.\\ 
Consider an observation $X_\epsilon \in \mx_\epsilon$ generated by a probability measure $\P_{\epsilon,\theta}$ on the measurable space $(\mx_\epsilon,\mathcal{A}_\epsilon)$ where the parameter $\theta$ is unknown. Given two families of measurable subsets $\Theta_{\epsilon,0},\Theta_{\epsilon,1}$, we will be testing the following hypotheses : 
\begin{equation}\label{hyptest}
    \H_0 : \theta \in \Theta_{\epsilon,0} \ \text{ vs } \  \H_1 : \theta \in \Theta_{\epsilon,1}
\end{equation}
and we will be studying the performance of tests as $\epsilon \rightarrow \epsilon_0$ .
The statistical performance of a family of tests $\{\psi_\epsilon\}$ can be studied using type I and II errors : \begin{equation*}
     \alpha_\epsilon(\psi_\epsilon) = \sup_{\theta \in \Theta_{\epsilon,0}}\E_{\epsilon,\theta} \psi_\epsilon \text{ and } \beta_\epsilon(\psi_\epsilon) = \sup_{\theta \in \Theta_{\epsilon,1}}\E_{\epsilon,\theta}[1-\psi_\epsilon] 
\end{equation*}
Or : \begin{equation*}
    \gamma_{\epsilon,t}(\psi_\epsilon) = t \alpha_\epsilon(\psi_\epsilon)+\beta_\epsilon(\psi_\epsilon) \ \text{ and } \ \gamma(t) = \inf_{\psi_\epsilon \in \Psi}\Bigl(t \alpha_\epsilon(\psi_\epsilon) + \beta_{\epsilon}(\psi_\epsilon)\Bigr)
\end{equation*}
In the minimax setting, we say that a family of tests $\psi_{\epsilon}^{(t)}$ is asymptotically minimax is : \begin{equation*}
    \gamma_{\epsilon,t}(\psi_{\epsilon}^{(t)}) = \gamma_\epsilon(t) + o(1) \ \text{ as } \  \epsilon \rightarrow \epsilon_0
\end{equation*}
Under the Neyman-Pearson setting, we consider families of tests $\psi_\epsilon$ satisfying $\alpha_\epsilon(\psi_\epsilon) \leq \alpha + o(1)$, and we say that a family $\psi_{\epsilon,\alpha}$ is asymptotically minimax if : 
\begin{equation*}
    \alpha_{\epsilon}(\psi_{\epsilon,\alpha}) \leq \alpha + o(1) \ \text{ and } \beta_\epsilon(\psi_{\epsilon,\alpha}) = \beta_\epsilon(\alpha) + o(1) \ \text{ as } \ \epsilon \rightarrow \epsilon_0
\end{equation*}
We are interested in the study of the asymptotic properties of $\gamma_\epsilon(t), \beta_\epsilon(\alpha)$, thanks to \Cref{conv}, this can be done only by studying the asymptotic properties of $\gamma_\epsilon \coloneqq \gamma_\epsilon(1)$. But this kind of problems is really difficult especially in the non parametric case, so we are interested 
in a simpler kind of problems called \textit{distinguishability problem} (\textit{rate asymptotic problem}) which can be summarized as follows : 
\begin{itemize}
    \item We want to obtain conditions for minimax distinguishability i.e. conditions for $\gamma_\epsilon \rightarrow 0$, and construct \textit{minimax consistent} families of tests $\psi_\epsilon$ such that $\gamma_\epsilon(\psi_\epsilon) \rightarrow 0$
    \item Obtain conditions for minimax nondistinguishability i.e. $\gamma_\epsilon \rightarrow 1$
\end{itemize}
\begin{prop}
    Let $\pi_{\epsilon,0},\pi_{\epsilon,1}$ be families of priors such that : 
    \begin{equation*}
        \pi_{\epsilon,0}(\Theta_{\epsilon,0}) \rightarrow 1 \ \text{ and } \ \pi_{\epsilon,1}(\Theta_{\epsilon,1}) \rightarrow 1 
    \end{equation*}
    Then, the following lower bounds hold true : 
    \begin{align*}
        \beta_\epsilon(\alpha) &\geq \beta_\epsilon(\alpha; \mix{\epsilon,0},\mix{\epsilon,1}) + o(1) \\
        \gamma_\epsilon &\geq \gamma_\epsilon(\mix{\epsilon,0},\mix{\epsilon,1}) + o(1) = 1 - \frac{1}{2}TV(\mix{\epsilon,0},\mix{\epsilon,1}) + o(1) 
    \end{align*}
\end{prop}
The problem that occurs often is that it is difficult to compute the total variation distance, so instead, we will use the $\chi^2$-divergence that one can define as follows : 
\begin{definition}
    Let $P_0,P_1$ be two probability measures on the same probability space. We define the $\chi^2$-divergence between $P_0$ and $P_1$ as follows : 
    \begin{equation}\label{chi2div}
        |P_0-P_1|_{2}^2 = \chi^2(P_0\Vert P_1) = \begin{cases*}
            \infty &\text{if} \ $P_0 \ \text{ does not dominate } \ P_1$ \\ 
            \E_{P_0} \bigg[\frac{dP_1}{dP_0}-1\bigg]^2 &\text{ otherwise }
        \end{cases*}
    \end{equation}
\end{definition}
Below, we will assume that $P_{\epsilon,0}$ dominates all the other measures $P_{\epsilon,\theta}$ for $\theta \in \Theta$ (This is fulfilled in the Gaussian sequence model)
\begin{prop}[$\chi^2$ nondistinguishability]\citet[Prop. 2.12]{Ing2012}
    \noindent
    \begin{itemize}
        \item Let $|P_{\epsilon,1}-P_{\epsilon,0}|_{2} \rightarrow 0$. Then $TV(P_{\epsilon,0},P_{\epsilon,1}) \rightarrow 0$ and $\gamma_\epsilon(P_{\epsilon,0},P_{\epsilon,1}) \rightarrow 0$.
        \item If $|P_{\epsilon,1}-P_{\epsilon,0}|_{2} = \bigo{1}$, then $\limsup TV(P_{\epsilon,0},P_{\epsilon,1}) < 2$ and $\liminf \gamma_\epsilon(P_{\epsilon,0},P_{\epsilon,1})>0$
    \end{itemize}
\end{prop}
\begin{proof}
The first statement follows from the inequalityc
\[
  |P_1 - P_0|_1 = E_{P_0}|L - 1| \leq \bigl(E_{P_0}|L - 1|^2\bigr)^{1/2} = |P_1 - P_0|_2.
\]
Let us establish the inequality
\begin{equation}\label{eq:2.58}
  |P_1 - P_0|_2 \geq |P_1 - P_0|_1 / 2\bigl(1 - \tfrac{1}{2}|P_1 - P_0|_1\bigr)^{1/2},
\end{equation}
which yields the second statement. Let us consider the optimal test $\psi = \psi^{(1)} =
\mathbf{1}_{\{L > 1\}}$ for testing $H_0 : P = P_0$ against $H_1 : P = P_1$. By~\eqref{eq:2.22}
we have
\[
  \gamma(\psi) = \gamma(P_0, P_1) = 1 - \tfrac{1}{2}|P_1 - P_0|_1 = 1 - E_{P_0}(L - 1)\psi.
\]
Then~\eqref{eq:2.58} follows from relations
\begin{align*}
  \tfrac{1}{2}|P_1 - P_0|_1
   &= E_{P_0}(L - 1)\psi \leq \bigl(E_{P_0}|L - 1|^2\bigr)^{1/2}\bigl(E_{P_0}\psi^2\bigr)^{1/2} \\
   &= |P_1 - P_0|_2\, \alpha^{1/2}(\psi) \leq |P_1 - P_0|_2\, \gamma^{1/2}(\psi) \\
   &= |P_1 - P_0|_2\bigl(1 - \tfrac{1}{2}|P_1 - P_0|_1\bigr)^{1/2}.
\end{align*}
\end{proof}
\subsection{Separation radius and goodness-of-fit}
In Goodness-of-Fit problems (i.e. $\Theta_{\epsilon,1} = \Theta_{\epsilon,0}^C$ in \Cref{hyptest}), it is impossible to use directly minimax approach because : 
\begin{equation*}
    \forall \epsilon : \gamma_\epsilon(\Theta_{\epsilon,0},\Theta_{\epsilon,1}) \equiv 1
\end{equation*}
Therefore, we consider alternatives of the form : 
\begin{equation*}
    \Theta_{\epsilon,1} = \bigl\{\theta \in \Theta : d(\theta,\Theta_{\epsilon,0}) \geq r_\epsilon\bigr\}
\end{equation*}
Where $d$ is a distance function and for $\theta \in \Theta : d(\theta,\Theta_{\epsilon,0}) \coloneqq \inf_{\theta_0 \in \Theta_{\epsilon,0}} d(\theta,\theta_0)$. Now, our problem can be stated as follows : What are the radii $r_\epsilon$ in order to obtain minimax distinguishability $\gamma_\epsilon \rightarrow 0$ ?\\ 
We will be studying this in the specific case of simple null-hypothesis of the form : $\H_0 : \theta \in \Theta_{\epsilon,0} = \{\theta_0\}$.
\begin{definition}[Critical radii]
We say that a family $r_{\epsilon}^*$ is the critical radii if it satisfies the following dichotomy : 
\begin{align*}
    \gamma_\epsilon \rightarrow 0 & \Leftrightarrow \frac{r_\epsilon}{r_{\epsilon}^*} \rightarrow \infty \\
    \gamma_\epsilon \rightarrow 1 &\Leftrightarrow \frac{r_\epsilon}{r_{\epsilon}^*} \rightarrow 0
\end{align*}
Moreover, we say that distinguishability conditions are in the classical form if the are in the form above. 
\end{definition}
\textcolor{green}{TO DO}
\subsection{Minimax properties of \texorpdfstring{$\chi^2$}{chi-square} type tests}
Now, return to the Gaussian model in the space of sequences $\l^2$, and condiser the following hypothesis problem : 
\begin{equation*}
    \H_0 : v = 0 \ \text{ versus } \ \H_1 : v \in V_\epsilon \subseteq \l^2
\end{equation*}
\begin{definition}[$\chi^2$-type tests]
    We say that a test $\psi$ is of $\chi^2$-type if there exists a family $w_\epsilon = \{w_{\epsilon,i}\}$ of sequences such that : \begin{equation*}
        \psi = \psi_{\epsilon,w_\epsilon;T} \coloneqq \ind{t_\epsilon > T_\epsilon} \ \text{ where } \ t_\epsilon = \sum_{i \in \N} w_{\epsilon,i}(x_i-1)^2
    \end{equation*}
\end{definition}
In the context of the definition above, set : \begin{equation*}
    h_\epsilon(v) = \sum_{i \in \N} w_{\epsilon,i}v_i^2; \ t_\epsilon(v) = t_\epsilon - h_\epsilon(v); \ |w_\epsilon| = \sup_{i \in \N} w_{\epsilon,i}
\end{equation*}
Under the null hypothesis, one can note that $\E_0 t_\epsilon$, moreover, one can assume the normalization $\sum w_{\epsilon,i}^2 = \frac{1}{2}$ in order to ensure that $\var{t_\epsilon}= 1$
\begin{lemma}[Minimax properties of $\chi^2$-type tests]\citet[Lemma 3.1 \& Corollary 3.3]{Ing2012}\\
    Assume that the family $w_\epsilon$ satisfies $w_{\epsilon,i} \geq 0$ as well as the normalization above. Then : \begin{equation}\label{chi2minimax1}
        \text{Var}_v(t_\epsilon) \leq 1 + 4 |w_\epsilon| h_\epsilon(v)
    \end{equation}
    Moreover, if $h_\epsilon \coloneqq \inf_{v \in V_\epsilon} h_\epsilon(v,w_\epsilon) \rightarrow \infty$ then for $T_\epsilon = \frac{h_\epsilon}{2}$, one has $\gamma(\psi_{w_\epsilon,T_\epsilon},V_\epsilon)$. 
    Additionally, if one assumes that $|w_\epsilon|$, then the statistics $t_\epsilon(v)$ are asymptotically $\mathbb{N}(0,1)$ uniformly over $v\in \l^2$ such that $h_\epsilon(v) = o(|w_\epsilon|^{-1})$, and the following inequalities hold true : 
    \begin{equation*}
        \beta(\psi_{\epsilon,T_\alpha},V_\epsilon) \leq \Phi(T_\alpha-h_\epsilon)+o(1), \ \gamma(\psi_{\epsilon,T_\epsilon}) \leq 2\Phi(-\frac{h_\epsilon}{2})+o(1)
    \end{equation*}
\end{lemma}
\begin{proof}
The proof combines Lemma~3.1 of Ingster--Suslina (giving the variance bound and asymptotic
normality of $t_\varepsilon$) with Corollary~3.3 (giving the minimax error bounds), the
latter itself drawing on Corollaries~3.1--3.2.

\medskip\noindent\textbf{Variance bound \eqref{eq:3.4}.}
Under $P_v$ we have $X_i = \eta_i + v_i$ with $\eta_i \sim N(0,1)$ i.i.d.\ Writing
$x = \eta + t$ with $\eta \sim N(0,1)$, $t \in \mathbb{R}$, a direct calculation gives
\begin{equation}\label{eq:3.11}
  E(x^2) = 1 + t^2, \quad \operatorname{Var}(x^2) = 2 + 4t^2, \quad
  E\bigl((x^2 - E x^2)^4\bigr) \leq c(1 + t^4),
\end{equation}
where $c$ is an absolute constant. Applying \eqref{eq:3.11} coordinatewise and using the
independence of the $X_i$,
\[
  E_v t_\varepsilon = \sum_i w_{\varepsilon,i}\, v_i^2 = h_\varepsilon(v),
  \qquad
  \operatorname{Var}_v(t_\varepsilon)
   = \sum_i w_{\varepsilon,i}^2\,(2 + 4 v_i^2)
   = 1 + 4\sum_i w_{\varepsilon,i}^2\, v_i^2,
\]
using the normalization $\sum_i w_{\varepsilon,i}^2 = 1/2$. Since
$w_{\varepsilon,i}^2 v_i^2 \leq |w_\varepsilon|\, w_{\varepsilon,i} v_i^2$, we get
\[
  \operatorname{Var}_v(t_\varepsilon)
   \;\leq\; 1 + 4|w_\varepsilon|\sum_i w_{\varepsilon,i}\, v_i^2
   \;=\; 1 + 4|w_\varepsilon|\, h_\varepsilon(v),
\]
which is~\eqref{eq:3.4}.

\medskip\noindent\textbf{Asymptotic normality of $t_\varepsilon(v)$.}
Set $\xi_{\varepsilon,i} = w_{\varepsilon,i}(X_i^2 - 1)$, so that $t_\varepsilon = \sum_i
\xi_{\varepsilon,i}$. By~\eqref{eq:3.11},
\[
  \sum_i E_v\bigl(\xi_{\varepsilon,i} - E_v\xi_{\varepsilon,i}\bigr)^4
   \;\leq\; c\sum_i w_{\varepsilon,i}^4\,(1 + v_i^4).
\]
Using $\sum_i w_{\varepsilon,i}^4 \leq |w_\varepsilon|^2 \sum_i w_{\varepsilon,i}^2 = |w_\varepsilon|^2/2$ and
$\sum_i w_{\varepsilon,i}^4 v_i^4 \leq \bigl(\sum_i w_{\varepsilon,i}^2 v_i^2\bigr)^2 \leq
|w_\varepsilon|^2 h_\varepsilon(v)^2$,
\[
  \sum_i E_v\bigl(\xi_{\varepsilon,i} - E_v\xi_{\varepsilon,i}\bigr)^4
   \;\leq\; c\,|w_\varepsilon|^2\bigl(1 + h_\varepsilon(v)^2\bigr).
\]
Since $\operatorname{Var}_v(t_\varepsilon) \geq 1$, the Lyapunov ratio is bounded by
$c|w_\varepsilon|^2(1 + h_\varepsilon(v)^2)$. Under the hypotheses
$|w_\varepsilon| = o(1)$ and $h_\varepsilon(v) = o(|w_\varepsilon|^{-1})$, this is $o(1)$
uniformly in the prescribed family, so the Lyapunov CLT yields
\[
  t_\varepsilon(v) = t_\varepsilon - h_\varepsilon(v)
   \;\xrightarrow{d}\; N\!\bigl(0,\,\operatorname{Var}_v(t_\varepsilon)\bigr)
   \;\xrightarrow{|w_\varepsilon| h_\varepsilon(v)\to 0}\; N(0, 1),
\]
uniformly over $v \in \ell^2$ such that $h_\varepsilon(v) = o(|w_\varepsilon|^{-1})$.

\medskip\noindent\textbf{Type I error and the bound $\gamma(\psi_{w_\varepsilon, T_\varepsilon}, V_\varepsilon) \to 0$.}
Under $H_0$ (i.e.\ $v = 0$) we have $h_\varepsilon(0) = 0$, so Step~2 gives
$t_\varepsilon \xrightarrow{d} N(0,1)$ under $P_0$. Hence $\alpha(\psi_{\varepsilon, T_\varepsilon}) =
\Phi(-T_\varepsilon) + o(1)$ for any deterministic threshold $T_\varepsilon$.

For Type II error in the consistency regime, fix $v \in V_\varepsilon$, so
$h_\varepsilon(v) \geq h_\varepsilon \to \infty$, and set $T_\varepsilon = h_\varepsilon/2$.
By Chebyshev,
\[
  \beta(\psi_{w_\varepsilon, T_\varepsilon}, v)
   = P_v(t_\varepsilon \leq T_\varepsilon)
   \leq P_v\bigl(|t_\varepsilon - h_\varepsilon(v)| \geq h_\varepsilon(v) - h_\varepsilon/2\bigr)
   \leq \frac{\operatorname{Var}_v(t_\varepsilon)}{(h_\varepsilon(v) - h_\varepsilon/2)^2}.
\]
Since $h_\varepsilon(v) \geq h_\varepsilon$ implies $h_\varepsilon(v) - h_\varepsilon/2 \geq
h_\varepsilon/2$, and $\operatorname{Var}_v(t_\varepsilon) \leq 1 + 4|w_\varepsilon|
h_\varepsilon(v)$,
\[
  \beta(\psi_{w_\varepsilon, T_\varepsilon}, v)
   \leq \frac{1 + 4|w_\varepsilon| h_\varepsilon(v)}{(h_\varepsilon/2)^2}
   = O\!\left(\frac{1}{h_\varepsilon^2} + \frac{|w_\varepsilon| h_\varepsilon(v)}{h_\varepsilon^2}\right) = o(1).
\]
Taking the supremum over $v \in V_\varepsilon$,
\[
  \gamma(\psi_{w_\varepsilon, T_\varepsilon}, V_\varepsilon)
   = \alpha(\psi_{w_\varepsilon, T_\varepsilon}) + \sup_{v \in V_\varepsilon}
       \beta(\psi_{w_\varepsilon, T_\varepsilon}, v) \to 0.
\]

\medskip\noindent\textbf{Sharp bounds under $|w_\varepsilon| = o(1)$.}
Now assume additionally $|w_\varepsilon| = o(1)$. By Step~3, it suffices to consider
$v \in V_\varepsilon$ with $h_\varepsilon(v) = O(1)$, since the case $h_\varepsilon(v) \to \infty$
yields $\beta \to 0$ already. For such $v$, Step~2 applies and gives
$t_\varepsilon(v) \xrightarrow{d} N(0, 1)$. Hence, with $T_\alpha = \Phi^{-1}(1 - \alpha)$,
\begin{align*}
  \beta(\psi_{\varepsilon, T_\alpha}, v)
   &= P_v\bigl(t_\varepsilon \leq T_\alpha\bigr)
    = P_v\bigl(t_\varepsilon - h_\varepsilon(v) \leq T_\alpha - h_\varepsilon(v)\bigr) \\
   &= \Phi\bigl(T_\alpha - h_\varepsilon(v)\bigr) + o(1)
    \leq \Phi(T_\alpha - h_\varepsilon) + o(1),
\end{align*}
using the monotonicity of $\Phi$ and $h_\varepsilon(v) \geq h_\varepsilon$. Taking the
supremum over $v \in V_\varepsilon$,
\[
  \beta(\psi_{\varepsilon, T_\alpha}, V_\varepsilon) \leq \Phi(T_\alpha - h_\varepsilon) + o(1).
\]

For $T_\varepsilon = h_\varepsilon/2$, the Type I error is
$\alpha(\psi_{\varepsilon, T_\varepsilon}) = \Phi(-h_\varepsilon/2) + o(1)$, while the
Type II error at any $v \in V_\varepsilon$ with $h_\varepsilon(v) = O(1)$ satisfies
\[
  \beta(\psi_{\varepsilon, T_\varepsilon}, v)
   = \Phi\bigl(h_\varepsilon/2 - h_\varepsilon(v)\bigr) + o(1)
   \leq \Phi(-h_\varepsilon/2) + o(1).
\]
Summing the Type I error and the supremum of the Type II error,
\[
  \gamma(\psi_{\varepsilon, T_\varepsilon})
   \leq \Phi(-h_\varepsilon/2) + \Phi(-h_\varepsilon/2) + o(1)
   = 2\Phi(-h_\varepsilon/2) + o(1). \qedhere
\]
\end{proof}
